Un film de savon vertical (indice \num{1.33}) entouré d'air est éclairé par une lumière de longueur d'onde \qty{\v{lam}}{nm} et observé en réflexion. Au début de l'observation, un point donné du film est au centre d'une frange brillante. À mesure que le film s'amincit, on compte \v{N} autres franges brillantes qui passent par ce point, puis le film y devient noir, juste avant d'éclater. L'observation a duré \qty{\v{T}}{s}.

\begin{enumerate}[label=\alph*)]
    \item De combien l'épaisseur du film change-t-elle entre deux franges brillantes successives ?
    \item Quelle était l'épaisseur du film en ce point au début de l'observation ?
    \item Quelle est la vitesse moyenne d'amincissement du film, en nanomètres par seconde ?
\end{enumerate}
